% 前言 (书页 v–viii / p005–p008)

\chapter*{前言}
\markboth{前言}{}
\addcontentsline{toc}{chapter}{前言}

知识的前沿（容我杜撰一个说法）总是在移动之中。今天的发现，明天就会成为每位科研工作者精神家什中的一件；到下周末，它就会出现在每一门研究生课程里；不出一个月，就会有人嚷着要把它写进本科课程表；再到明年，我确信，它将显得如此寻常，以至于可以假定连每个学童都早已知道。

把拓居的边界向前推进、垦殖并教化新领地的过程，是分阶段进行的。首先是原始论文的发表：作者们满心欢喜，读者们冷眼审视。随后是综述文章，提供粗略的草图规划——穿越文献丛林的初阶向导。再往后是各部专著，如同精确的地籍测量，把已经夺得的疆域绘制成图，调停优先权之争，让每一件事实、每一个理论各归其位。

最后，我们需要教科书。专著（treatise）与教科书（textbook）之间有一条深刻的分界：专著重阐说；教科书重讲解。从来没有人设想学生能把整个物理学装进脑袋，甚至连物理学任何一个分支的全部也装不进去。凡是翻检专著、综述与原始论文便唾手可得的东西，他无须记诵；但他必须学会\emph{阅读}那些文献：必须学会它们所用的语言；必须知道他的这门科学赖以奠基的基本实验事实和一般理论原理。

本书的目标，是以尽可能简单的方式呈示完美晶态固体物理学理论的要素。这是一本充满\emph{观念}（ideas）而非事实的书；它是对原理的阐述，而不是对现象的描述。

理论是对某个假想模型的性质所作的分析。物理学几乎可以定义为把宇宙归并于数学的智识操练；在这门学科里，我们的模型是数学的。本书讨论的理论都是数学理论；这个领域中最重要的概念，比如“费米面（Fermi Surface）”，都是抽象的数学构造，不经形式化的分析，便无法得到恰当的解释和理解。

我所力图做到的，是对能够体现每一条原理的最简单模型，给出自足的数学处理。如果超导体的大多数有趣性质都能从一个带有削减吸引相互作用的自由电子模型推导出来，那么这就是计算的框架。在这个阶段，与其为了预见观测到的对简单公式的种种偏离，而挥舞那些基于更真实却复杂得多的物理设定、因而远为笨重的方程，倒不如先领会该现象之所以能够出现所必需的条件。至于初等模型的这些精细发挥，读者必须去求诸原始论文和专著。

另一方面，模型既经界定，就不可在数学分析面前畏缩。以我的经验，许多简单而著名的公式，要从第一性原理出发直接推导，在印刷物中并不容易找到：原始论文走的不是最便当的路子；综述的作者们觉得必要的推演太难——或者有失身份；专著又对完备与严格过分矜持。我力图使数学论证自身完整——至少在原理上可以理解——而不必动辄祈灵于疲惫的作者的那尊 \emph{deus ex machina}（解围之神）——“可以证明……”。试图覆盖如此广阔的领域有一个好处：可以援引普遍原理，例如 Bloch 定理与区理论（the theory of zones），把这门学科的许多分支统一起来，省下许多脑力。

那么预期读者已经懂得多少呢？他应当熟悉关于固体的那些初等的描述性事实——例如金属的自由电子理论——即本科课程所讲授的内容。我还假定他具备量子力学的初步知识——特别是 Schrödinger 方程、微扰论和散射理论——正如如今的实验物理学研究生被期望掌握的那样。我设法把数学技巧保持在这个水平上；每当代数眼看要变得艰难，我就停了下来。书中没有密度矩阵、气泡图、分支点、特征标表，也没有高级理论武库中的其他零碎什物；职业的“理论家”们必须另觅粮草。

为着那些以一本书缺了什么来评判它的书评家着想，让我承认：书中没有对合金、位错、F 心（F-centre）、杂质中心等等的认真讨论；与单个原子或原子核相联系的磁共振一字不提；也无意去解释本质上属于宏观的现象，如铁磁畴（ferromagnetic domains）、$p$--$n$ 结或超导体的中间态。把“简单、完美、晶态固体”以外的一切统统排除在外，这种做法是人为的，\footnote{译注：原文作 “the exclusion of all 'except simple, perfect, crystalline solids”，其中 “'except” 是 except 的口语省音拼法，义即“除了”。}但它很便利，因为它赋予论述某种统一性：全书的中心，是晶格周期性的数学推论。再者，朋友们，人生苦短。

对物理学这样一个活跃分支的这番讲述，并不自命有何新意。这些理论正是当今在这个领域工作的人们所实际使用、并承认为充分确立的理论。本书不去批判这些理论，不去讨论它们作为自然现象之解释的有效性，不以充分的数学严格性推导它们，也不展示其应用的繁花盛果。全部心力都放在表述的清晰上。我们要求读者掌握新概念；那就请他暂且悬置自己的批判才能，直到他理解了这些新观念的含义；到那时，倘若他合乎情理地依然怀疑，就让他转向这门科学赖以建成的浩瀚文献，从那些虽浑浊却丰沛的源泉中，为他的“不信”寻求援助。我有意避免对原始论文作任何直接引用；需要这类信息的学生，请查阅书末列出的专著和综述文章。

本书最初是剑桥大学理论与实验物理学研究生课程的一份讲义：它是在我仍是 Cavendish 实验室在编人员的最后两个学期里写成的，并在离任前的最后几个星期里完成。九年之间，得隶身于那座无与伦比的机构，是极大的荣幸。我也再清楚不过：要在物理学世界里企及它所树立、并且还在继续树立的那个独一无二的标准，是何等不可能的事。

但是 Cavendish 不只是一座著名的实验室；它更是人性与友谊的居所。请允许我向那些朋友致谢——尤其是 Nevill Mott、Brian Pippard 与 Volker Heine——他们把我带到剑桥；他们迎纳我、教导我、有见地地驳难我、不遗余力地扶助我，使我在此间的生活惬意、有趣而令人振奋。他们聆听过天球的音乐；然而他们知道，科学是为人而设的，人却不是为科学而设。

\begin{flushright}
J. M. Z.
\end{flushright}

\noindent\emph{Cavendish Laboratory, Cambridge}\\
1963 年 6 月

\chapter*{第二版前言}
\markboth{第二版前言}{}
\addcontentsline{toc}{chapter}{第二版前言}

这个新版本仍然打算遵循上面那篇前言所申明的诸原则。但是八年大约正是现代科学知识的倍增时间，而这期间固态物理学家们并没有闲着。原文的大部分依然成立，但增加了若干新的节，用以覆盖近来变得更加引人注目的论题，或者在理解上、侧重点上出现了显著转变的论题。我还试图顺带提及一些与基础理论相关的现象或研究领域，即使无法详细讨论。这样，本书的总体范围有所拓宽，例如纳入了关于磁性与非磁性杂质、F 心、表面、隧穿、结，以及第二类超导电性（type II superconductivity）的若干内容。但是数学精致程度的总体水平并没有提高，尽管高级量子理论的技术形式体系如今正在这个领域里变得愈来愈常见。

我深深感激许多同行——尤其是 Bristol 的 Bob Chambers 与马德里的 Federico Garcia-Moliner——感谢他们提出的大量细致意见；我在新文本中已尽力加以采纳。Bob Evans 编制了一部新索引，帮了大忙。为了免得读者觉得离开剑桥是一段漫长的流放岁月，请容我只补一句：Bristol 同样是个毫不逊色的、值得投奔的“好窟”。\footnote{译注：原文作 “Bristol, too, is just as good a 'ole to go to”。其中 “'ole” 是 “hole”（洞、窟）的伦敦俚音拼法；作者以自嘲的口吻把自己的新工作地点称作一个“窟”——然而它是与剑桥同样好的“窟”。}

\begin{flushright}
J. M. Z.
\end{flushright}

\noindent\emph{H. H. Wills Physics Laboratory, Bristol}\\
1971 年 3 月
